\documentclass[11pt]{article}
\usepackage{mathblatt}
\makeatletter
\setlength{\mbpfbe}{0mm}% keine Punkte (Bauregel 6.4)
\makeatother
% Rechenraum als Karo (Bauregel 6.9): n Rechenschritte = n mal 10 mm
\newcommand{\karo}[1]{\par\smallskip\noindent\begin{tikzpicture}
  \clip (0,0) rectangle ({\linewidth-1mm},{#1*10mm});
  \draw[black!16,line width=0.3pt,step=5mm] (0,0) grid ({\linewidth},{#1*10mm});
  \end{tikzpicture}\par}
\newcommand{\linie}[1]{\,{\color{black!45}\rule[-1pt]{#1}{0.4pt}}\,}
\newcommand{\pkt}{\textperiodcentered{}}
\begin{document}

\pfheftstil
\pflaufkopf[P10 \pkt{} mehr]{Lineare Funktionen}{PA10-6}\pfseitevon
\pfheftkopf{Lineare Funktionen}{mehr}
\pfstufenzeile{Punktprobe \pkt{} Gleichung aus Text \pkt{} Gerade zeichnen \pkt{} Nullstelle \pkt{} Gleichung aus zwei Punkten}
% 2017-OS-K5b-m1 (Schwere I, rang 4)
\begin{pfaufg}{nach~P10\\’17~\pkt{}~5b}{1.}{}
Die Gerade $g$ hat die Gleichung $y = \tfrac{2}{3}x + 2$. Auf ihr liegt der Punkt $A(-12\,|\,y)$. Berechne seine y-Koordinate.
\karo{2}
\fusshilfe{\mbox{1:~$-$6}}
\end{pfaufg}
% 2021-OS-K6a-m1 (Schwere I, rang 5)
\begin{pfaufg}{nach~P10\\’21~\pkt{}~6a}{2.}{}
Eine 36 cm hohe Kerze wird angezündet. Die Tabelle hält fest, wie hoch sie nach einigen Minuten noch ist. Trag die zwei fehlenden Werte in die Tabelle ein.
\par\smallskip\noindent{\small\itshape Skizze: Tabelle: Zeit in Minuten | 0 | 10 | 20 | 30 | 40 | … | 80; Höhe der Kerze in cm | leer | 33 | 30 | 27 | 24 | … | leer}\par
\karo{3}
\fusshilfe{\mbox{2:~0 min: 36 cm; 80 min: 12 cm}}
\end{pfaufg}
% 2021-OS-K6d-m1 (Schwere I, rang 6)
\begin{pfaufg}{nach~P10\\’21~\pkt{}~6d}{3.}{}
Eine 36 cm hohe Kerze wird angezündet. Die Tabelle hält fest, wie hoch sie nach einigen Minuten noch ist. Die Gleichung $y = -0{,}3x + 36$ beschreibt das Abbrennen der Kerze ($x$: Zeit in Minuten, $y$: Höhe in cm). Bestimme, nach wie vielen Minuten die Kerze ganz heruntergebrannt ist.
\par\smallskip\noindent{\small\itshape Skizze: Tabelle: Zeit in Minuten | 0 | 10 | 20 | 30 | 40 | … | 80; Höhe der Kerze in cm | leer | 33 | 30 | 27 | 24 | … | leer}\par
\karo{2}
\fusshilfe{\mbox{3:~120 min}}
\end{pfaufg}
% 2021-OS-K7a-m1 (Schwere I, rang 7)
\begin{pfaufg}{nach~P10\\’21~\pkt{}~7a}{4.}{}
Ordne jedem Sachverhalt die passende Gleichung zu. Sachverhalt 1: Eine Taxifahrt kostet 4 € Grundgebühr und 30 Cent je Kilometer. Sachverhalt 2: Ein Fitnessstudio verlangt im Monat 15 € Grundgebühr und 4 € für jeden gebuchten Kurs. Gleichungen: $y = 30x + 4$, $y = 4x + 15$, $y = 0{,}30x + 4$.
\par\smallskip\noindent{\small\itshape Skizze: Ankreuztabelle: Zeilen Sachverhalt 1, Sachverhalt 2; Spalten $y = 30x + 4$ | $y = 4x + 15$ | $y = 0{,}30x + 4$}\par
\karo{3}
\fusshilfe{\mbox{4:~4x + 15}}
\end{pfaufg}
% 2016-OS-K6a-m1 (Schwere I, rang 8)
\begin{pfaufg}{nach~P10\\’16~\pkt{}~6a}{5.}{}
Die 10b fährt mit dem Bus in eine Jugendherberge. Zwei Busfirmen machen ein Angebot: „Sonnenschein“ verlangt 150 € Grundpreis und 1,20 € je gefahrenem Kilometer, „Reiselust“ keinen Grundpreis und 1,80 € je Kilometer. Die Grafik zeigt für beide Angebote die Fahrtkosten in Abhängigkeit von der Strecke. Gib zu jedem der beiden Graphen die passende Busfirma an. Begründe für „Sonnenschein“, warum dein Graph passt.
\par\smallskip\noindent{\small\itshape Skizze: Tabelle: Busfirma | Grundpreis | Preis je km; Sonnenschein | 150 € | 1,20 €; Reiselust | – | 1,80 € // Graph: I: 1.8*x | II: 1.2*x+150; x 0..350; y 0..700; Gitter ja; Achsen Fahrtstrecke in km | Fahrtkosten in €}\par
\karo{3}
\fusshilfe{\mbox{5:~II beginnt bei 150 € (Grundpreis bei 0 km) und steigt flacher (1,20 €/km)}}
\end{pfaufg}
% 2021-OS-K2a-m1 (Schwere I, rang 8)
\begin{pfaufg}{nach~P10\\’21~\pkt{}~2a}{6.}{}
Die lineare Funktion $f$ hat die Gleichung $y = 2x - 1$. Zeichne den Graphen von $f$ in das Koordinatensystem.
\par\smallskip\noindent{\small\itshape Skizze: Koordinatensystem: x -4..4; y -4..4; Gitter ja}\par
\karo{2}
\fusshilfe{\mbox{6:~Gerade durch (0|$-$1) und (1|1)}}
\end{pfaufg}
% 2021-OS-K2b-m1 (Schwere I, rang 9)
\begin{pfaufg}{nach~P10\\’21~\pkt{}~2b}{7.}{}
Die lineare Funktion $f$ hat die Gleichung $y = 2x - 1$. Kreuze die beiden Aussagen an, die für $f$ stimmen:
\par\smallskip\noindent{\small\itshape Skizze: Ankreuztabelle: vier Aussagen, je ein Ankreuzfeld}\par
\par\smallskip\noindent\pfkreuz{A}{Der Punkt $Q(4\,|\,0)$ liegt auf dem Graphen.}\pfkreuz{B}{Der Graph steigt monoton.}\pfkreuz{C}{Der Graph trifft die y-Achse in $P(0\,|\,-1)$.}\pfkreuz{D}{Der Graph geht durch den Ursprung.}
\fusshilfe{\mbox{7:~Aussage 4 falsch}}
\end{pfaufg}
% 2021-OS-K6c-m1 (Schwere II, rang 9)
\begin{pfaufg}{nach~P10\\’21~\pkt{}~6c}{8.}{}
Eine 36 cm hohe Kerze wird angezündet. Die Tabelle hält fest, wie hoch sie nach einigen Minuten noch ist. Das Abbrennen der Kerze beschreibt die Gleichung $y = -0{,}3x + 36$. Schreib zu $y$, $x$ und $36$, was sie hier bedeuten. Beispiel: $-0{,}3$ heißt, dass die Kerze in jeder Minute um 0,3 cm kürzer wird.
\par\smallskip\noindent{\small\itshape Skizze: Tabelle: Teil der Gleichung | Bedeutung; y | leer; $-$0,3 | In jeder Minute wird die Kerze 0,3 cm kürzer.; x | leer; 36 | leer}\par
\karo{4}
\fusshilfe{\mbox{8:~36: Anfangshöhe der Kerze in cm}}
\end{pfaufg}
% 2016-OS-K6b-m1 (Schwere II, rang 10)
\begin{pfaufg}{nach~P10\\’16~\pkt{}~6b}{9.}{}
Die 10b fährt mit dem Bus in eine Jugendherberge. Zwei Busfirmen machen ein Angebot: „Sonnenschein“ verlangt 150 € Grundpreis und 1,20 € je gefahrenem Kilometer, „Reiselust“ keinen Grundpreis und 1,80 € je Kilometer. Die Grafik zeigt für beide Angebote die Fahrtkosten in Abhängigkeit von der Strecke. Die Jugendherberge liegt 110 km entfernt; der Bus fährt hin und zurück. Gib an, welche Busfirma dafür billiger ist. Am Ort kommt noch ein Tagesausflug von 60 km hinzu. Untersuche, ob dieselbe Firma dann noch die billigere ist, und begründe.
\par\smallskip\noindent{\small\itshape Skizze: Tabelle: Busfirma | Grundpreis | Preis je km; Sonnenschein | 150 € | 1,20 €; Reiselust | – | 1,80 € // Graph: I: 1.8*x | II: 1.2*x+150; x 0..350; y 0..700; Gitter ja; Achsen Fahrtstrecke in km | Fahrtkosten in €}\par
\karo{3}
\fusshilfe{\mbox{9:~Reiselust billiger: Sonnenschein 414 € gegen Reiselust 396; nein: bei 280 km Reiselust 504 €, Sonnenschein 486}}
\end{pfaufg}
% 2019-OS-K2a-m1 (Schwere I, rang 10)
\begin{pfaufg}{nach~P10\\’19~\pkt{}~2a}{10.}{}
In einem Koordinatensystem sind die Graphen $f$, $g$ und $h$ von drei linearen Funktionen gezeichnet. Gib die Nullstelle der Funktion $g$ an.
\par\smallskip\noindent{\small\itshape Skizze: Graph: f: -3*x+3 | g: x+3 | h: -1; x -5..5; y -5..5; Gitter ja}\par
\karo{2}
\fusshilfe{\mbox{10:~$-$3}}
\end{pfaufg}
% 2016-OS-K6c-m1 (Schwere II, rang 11)
\begin{pfaufg}{nach~P10\\’16~\pkt{}~6c}{11.}{}
Die 10b fährt mit dem Bus in eine Jugendherberge. Zwei Busfirmen machen ein Angebot: „Sonnenschein“ verlangt 150 € Grundpreis und 1,20 € je gefahrenem Kilometer, „Reiselust“ keinen Grundpreis und 1,80 € je Kilometer. Die Grafik zeigt für beide Angebote die Fahrtkosten in Abhängigkeit von der Strecke. Gib die Gleichung der Funktion an, deren Graph II ist (x: Strecke in km, y: Kosten in €).
\par\smallskip\noindent{\small\itshape Skizze: Tabelle: Busfirma | Grundpreis | Preis je km; Sonnenschein | 150 € | 1,20 €; Reiselust | – | 1,80 € // Graph: I: 1.8*x | II: 1.2*x+150; x 0..350; y 0..700; Gitter ja; Achsen Fahrtstrecke in km | Fahrtkosten in €}\par
\karo{2}
\fusshilfe{\mbox{11:~1,2x + 150 (x in km, y in €)}}
\end{pfaufg}
% 2016-OS-B1c-m1 (Schwere I, rang 12)
\begin{pfaufg}{nach~P10\\’16~\pkt{}~1c}{12.}{}
Im Koordinatensystem sind zwei Geraden $f$ und $g$ gezeichnet. Kreuze an, welche der beiden fällt:
\par\smallskip\noindent{\small\itshape Skizze: Graph: f: -1*x+1 | g: 0.75*x-1.5; x -1.5..3.5; y -2.5..1.5; Gitter ja}\par
\par\smallskip\noindent\pfkreuz{A}{$f$}\pfkreuz{B}{$g$.}
\fusshilfe{\mbox{12:~f}}
\end{pfaufg}
% 2017-OS-K5a-m1 (Schwere II, rang 14)
\begin{pfaufg}{nach~P10\\’17~\pkt{}~5a}{13.}{}
Eine Gerade $g$ geht durch die Punkte $K(-6\,|\,-2)$ und $L(3\,|\,4)$. Zeichne ein Koordinatensystem und darin die Gerade $g$. Zeige, dass $y = \tfrac{2}{3}x + 2$ eine Gleichung von $g$ ist.
\par\smallskip\noindent{\small\itshape Skizze: leeres Karoraster ohne Achsen (etwa 28 × 15 Kästchen)}\par
\karo{4}
\fusshilfe{\mbox{13:~Gerade durch K und L (y-Achsenabschnitt 2, Nullstelle $-$3); $\tfrac{2}{3}$x + 2}}
\end{pfaufg}
% 2015-OS-K4d-m1 (Schwere II, rang 15)
\begin{pfaufg}{nach~P10\\’15~\pkt{}~4d}{14.}{}
Lena springt mit dem Fallschirm aus einem Flugzeug. Ihre Höhe $h$ (in Metern) hängt von der Zeit $t$ (in Sekunden) ab. Nach 25 Sekunden freiem Fall öffnet sie in 1500 m Höhe den Schirm; 100 Sekunden später ist sie noch 900 m hoch. Während sie mit offenem Schirm gleitet, liegt der Graph auf einer Geraden. Stelle eine Gleichung dieser Geraden auf.
\par\smallskip\noindent{\small\itshape Skizze: Graph: G: -2*x$^{2}$+2750 für 0..25 | G: -6*x+1650 für 25..275; x 0..300; y 0..3000; Gitter ja; Achsen Zeit t in Sekunden | Höhe h in Metern}\par
\karo{2}
\fusshilfe{\mbox{14:~$-$6t + 1650}}
\end{pfaufg}
\par\vfill\noindent{\scriptsize\color{mbgrau}Originale: P10 2017 \pkt{} 5b, P10 2021 \pkt{} 6a, P10 2021 \pkt{} 6d, P10 2021 \pkt{} 7a, P10 2016 \pkt{} 6a, P10 2021 \pkt{} 2a, P10 2021 \pkt{} 2b, P10 2021 \pkt{} 6c, P10 2016 \pkt{} 6b, P10 2019 \pkt{} 2a, P10 2016 \pkt{} 6c, P10 2016 \pkt{} 1c, P10 2017 \pkt{} 5a, P10 2015 \pkt{} 4d\par}
\end{document}